\begin{frame}
Covariance Analysis: Overview
\begin{itemize}
  \item Covariance propagation and covariance update can be be accomplished without actually dealing with the state vector at all.
\begin{equation}\label{altcovextrapeq}
\altcovextrapeq
\end{equation}
\begin{equation}\label{covupdateeq}
  \covupdateeq
\end{equation}
  \item Of course, the Kalman gain still has to be computed for use in Eq. \eqref{covupdateeq} above, the update equation.  The gain is:
\begin{equation}\label{kalgaineq}
\kalgaineq
\end{equation}
\end{itemize} 
\end{frame}

\begin{frame}
Covariance Analysis: Overview
\begin{itemize}
  \item As in real-time implementation, the Joseph form of the gain equation is generally preferred, because of its superior numerical properties:
\begin{equation}
\josephcovupdateeq
\end{equation} 
\end{itemize} 
\end{frame}

\begin{frame}
Covariance Analysis: Overview
\begin{itemize}
  \item The only items driving the covariance equations are the state dynamics $F$, the process noise $Q$, and the measurement noise $R$
  \item Thus, covariance analysis can provide  good insight into how well balanced the $Q$ and $R$ matrices are
  \item More importantly, covariance analysis is crucial in exploring what-if scenarios with new measurement sources
\end{itemize} 
\end{frame}

\begin{frame}
Covariance Analysis: Overview
\begin{itemize}
  \item As an example, suppose you want to evaluate several different GPS receivers as navaids
  \item You can use a manufacturer's published $1\sigma$ accuracy specifications in your measurement noise matrix $R$ and run several covariance simulations to see what the expected performance is
\end{itemize} 
\end{frame}

\begin{frame}
Covariance Analysis
\begin{itemize}
  \item Of course, if GPS is your only navaid, this is a no brainer:  the receiver with the least measurement noise performs best
  \item However, if this sensor is just one of several, you might find that some less expensive receiver when thrown into the mix with other navaids is good enough 
\end{itemize} 
\end{frame}

\begin{frame}
Covariance Analysis: Overview
\begin{itemize}
  \item Other uses: Debugging the filter
  \item If you save the actually used process noise $Q$ and measurement noise $R$ from a real-time mission (using a maintenance logging facility), then thse matrices can be fed into the covariance simulation to see what effect these have on the gains.
  \item Or to see how the offline gains compare with the real-time gains.
\end{itemize} 
\end{frame}

\begin{frame}
Covariance Analysis: Overview
\begin{itemize}
  \item Covariance analysis is quite often used to determine proper error budgets
  \item Error budget:  How error sources are allocated to various states and sensors
  \item E.g., in a nav application, to meet overall system position accuracy requirements, certain combinations of misalignement, instrument bias, scale factor, and inertial noise error, etc., will work when coupled with sensor errors
\end{itemize} 
\end{frame}

\begin{frame}
Covariance Analysis: Overview
\begin{itemize}
  \item Optimal vs. Suboptimal
  \item Evaluating finite wordlength effects 
  \item Bottom Line: Covariance analysis essential to filter maintenance
	\begin{itemize}
		\item If you are planning to modify a Kalman filter at all, you need to run covariance analysis simulations
		\item It is very risky not to do so!
	\end{itemize}
\end{itemize} 
\end{frame}
